\section{因果性}\label{sec:causality-in-lorentzian-manifolds}

本节研究时空的因果性质. Riemann 几何中没有对应的课题, 因为那里的度量是正定的. 虽然我们将主要讨论四维 Lorentz 流形, 讨论也容易推广到任意维数 $n\geq2$.

若存在向量场 $X\in\mathfrak{X}(M)$ 使 $\langle X,X\rangle<0$, 则称时空 $(M,g)$ \textbf{时间可定向}. 此时, 可以把连通分支 $C(X_p)$ 中的类时向量选为指向未来的向量, 从而在每个切空间 $T_pM$ 上定义时间定向 (每个切空间当然都与 Minkowski 时空等距).

假定 $(M,g)$ \textbf{已作时间定向} (即时间可定向, 且已选定时间定向). 若 $\dot c$ 指向未来, 则称类时曲线 $c:I\subset\mathbb{R}\to M$ \textbf{朝向未来}. $p\in M$ 的\textbf{时序未来} $I^+(p)$ 是可以由朝向未来的类时曲线从 $p$ 到达的所有点所成的集合. \textbf{朝向未来的因果曲线}是满足 $\dot c$ 为类时或零模且 (若非零) 指向未来的曲线 $c:I\subset\mathbb{R}\to M$. $p\in M$ 的\textbf{因果未来} $J^+(p)$ 是可以由朝向未来的因果曲线从 $p$ 到达的所有点所成的集合. 注意, $I^+(p)$ 就是从 $p$ 出发的非零质量粒子能到达的所有事件所成的集合, 而 $J^+(p)$ 是可以被 $p$ 因果地影响的事件所成的集合 (因为这种因果影响不能传播得比光更快). 类似地, $p\in M$ 的\textbf{时序过去} $I^-(p)$ 是可以由朝向未来的类时曲线到达 $p$ 的所有点所成的集合; $p\in M$ 的\textbf{因果过去} $J^-(p)$ 是可以由朝向未来的因果曲线到达 $p$ 的所有点所成的集合.

一般而言, 由于时空的整体性质, 时序过去与未来以及因果过去与未来可为相当复杂的集合. 然而在局部, 因果性质与 Minkowski 时空中的类似. 更确切地说, 有下列论断:

\begin{RiemannProposition}\label{prop:local-causal-structure-in-convex-normal-neighborhood}
设 $(M,g)$ 为已作时间定向的时空. 则每点 $p_0\in M$ 都有一个开邻域 $V\subset M$, 使将 $g$ 限制到 $V$ 所得的时空 $(V,g)$ 满足:
\begin{enumerate}[label=(\arabic*),ref=\arabic*,leftmargin=*]
\item\label{cond:geodesically-convex-normal-neighborhood}
$V$ 是其每一点的正规邻域, 且对任意 $p,q\in V$, 都存在唯一的 (在重新参数化意义下) 连接 $p,q$ 的测地线 (即 $V$ 是\textbf{测地凸的});
\item\label{cond:local-chronological-future-via-timelike-geodesics}
$q\in I^+(p)$, 当且仅当存在连接 $p$ 至 $q$ 的朝向未来的类时测地线;
\item\label{cond:local-causal-future-is-chronological-closure}
$J^+(p)=\overline{I^+(p)}$;
\item\label{cond:local-boundary-via-null-geodesics}
$q\in J^+(p)\setminus I^+(p)$, 当且仅当存在连接 $p$ 至 $q$ 的朝向未来的零模测地线.
\end{enumerate}
\end{RiemannProposition}

\begin{RiemannProof}
设 $U$ 为 $p_0$ 的正规邻域, 在 $U$ 上选取由参数化
\[
  \varphi(x^0,x^1,x^2,x^3)=\exp_{p_0}(x^0v_0+x^1v_1+x^2v_2+x^3v_3)
\]
给出的\textbf{正规坐标} $(x^0,x^1,x^2,x^3)$, 其中 $\{v_0,v_1,v_2,v_3\}$ 是 $T_{p_0}(M)$ 的一组基 [参见第~\ref{chap:riemannian-manifolds}~章的 Exercise~\ref{ex:minimizing-properties-of-geodesics}(\ref{ex:metric-and-connection-in-normal-coordinates})]. 设 $D:U\to\mathbb{R}$ 为可微函数
\[
  D(p):=\sum_{\alpha=0}^3(x^\alpha(p))^2,
\]
并对每个 $\varepsilon>0$ 定义集合
\[
  B_\varepsilon:=\{p\in U\mid D(p)<\varepsilon\}.
\]
对于足够小的 $\varepsilon$, 该集合与 $T_{p_0}M$ 中的一个开球微分同胚. 为简单起见, 假定 $U$ 就具有这个形式.

下面证明, 存在 $\rho>0$, 使当 $c:I\subset\mathbb{R}\to B_\rho$ 是测地线时, $D(t):=D(c(t))$ 的所有临界点都是严格局部极小点. 事实上, 令 $x^\mu(t):=x^\mu(c(t))$, 有
\begin{align*}
  \dot D(t)&=2\sum_{\alpha=0}^3 x^\alpha(t)\dot x^\alpha(t);\\
  \ddot D(t)&=2\sum_{\alpha=0}^3(\dot x^\alpha(t))^2
       +2\sum_{\alpha=0}^3x^\alpha(t)\ddot x^\alpha(t)\\
       &=2\sum_{\mu,\nu=0}^3\left(\delta_{\mu\nu}
            -\sum_{\alpha=0}^3\Gamma_{\mu\nu}^\alpha(c(t))x^\alpha(t)\right)
            \dot x^\mu(t)\dot x^\nu(t),
\end{align*}
且对足够小的 $\rho$, 矩阵
\[
  \delta_{\mu\nu}-\sum_{\alpha=0}^3\Gamma_{\mu\nu}^\alpha x^\alpha
\]
在 $B_\rho$ 上正定.

考虑由
\[
  F(v)=(\pi(v),\exp(v))
\]
定义的映射 $F:W\subset TM\to M\times M$, 其中 $W$ 为 $0\in T_{p_0}M$ 的某个开邻域. 如同在 Riemann 情形中所见 (参见第~\ref{chap:riemannian-manifolds}~章第~\ref{sec:minimizing-properties-of-geodesics}~节), 此映射在 $0\in T_{p_0}M$ 处是局部微分同胚. 取足够小的 $\delta>0$ 并缩小 $W$, 可假定 $F$ 将 $W$ 微分同胚地映到 $B_\delta\times B_\delta$, 且对任意 $t\in[0,1]$, $v\in W$ 都有 $\exp(tv)\in B_\rho$ (否则, 就能构造一个满足 $v_n\to0\in T_{p_0}M$ 而 $\exp(v_n)\not\to p_0$ 的序列).

最后, 令 $V=B_\delta$. 若 $p,q\in V$, $v=F^{-1}(p,q)$, 则 $c(t)=\exp_p(tv)$ 是连接 $p$ 与 $q$ 的测地线, 其像包含在 $B_\rho$ 中. 若其像没有包含在 $V$ 中, 则 $D(t)$ 必有一个局部极大点, 这是不可能的. 因此在 $V$ 中有一条测地线连接 $p$ 与 $q$. 由于 $\exp_p$ 是到 $V$ 的微分同胚, 测地线是唯一的 (不计重新参数化). 这证明了 (\ref{cond:geodesically-convex-normal-neighborhood}).

为证明论断 (\ref{cond:local-chronological-future-via-timelike-geodesics}), 先注意若有朝向未来的类时测地线连接 $p$ 至 $q$, 则显然 $q\in I^+(p)$. 现假设 $q\in I^+(p)$; 则存在朝向未来的类时曲线 $c:[0,1]\to V$, 使 $c(0)=p$, $c(1)=q$. 选取由参数化
\[
  \varphi(x^0,x^1,x^2,x^3)=\exp_p(x^0E_0+x^1E_1+x^2E_2+x^3E_3)
\]
给出的正规坐标 $(x^0,x^1,x^2,x^3)$, 其中 $\{E_0,E_1,E_2,E_3\}$ 是 $T_pM$ 的标准正交基, $E_0$ 为指向未来的类时向量. 由于 $F:W\to V\times V$ 是微分同胚, 这些是 $V$ 中的全局坐标. 定义
\begin{align*}
  W_p(q)&=-(x^0(q))^2+(x^1(q))^2+(x^2(q))^2+(x^3(q))^2\\
        &=\sum_{\mu,\nu=0}^3\eta_{\mu\nu}x^\mu(q)x^\nu(q),
\end{align*}
其中 $(\eta_{\mu\nu})=\operatorname{diag}(-1,1,1,1)$, 则需要证明 $W_p(q)<0$. 令 $W_p(t):=W_p(c(t))$. 由于 $x^\mu(p)=0$ ($\mu=0,1,2,3$), 有 $W_p(0)=0$. 令 $x^\mu(t):=x^\mu(c(t))$, 得到
\begin{align*}
  \dot W_p(t)&=2\sum_{\mu,\nu=0}^3\eta_{\mu\nu}x^\mu(t)\dot x^\nu(t);\\
  \ddot W_p(t)&=2\sum_{\mu,\nu=0}^3\eta_{\mu\nu}x^\mu(t)\ddot x^\nu(t)
               +2\sum_{\mu,\nu=0}^3\eta_{\mu\nu}\dot x^\mu(t)\dot x^\nu(t).
\end{align*}
因而 (回顾 $(\mathrm{d}\exp_p)_0=\operatorname{id}$)
\begin{align*}
  \dot W_p(0)&=0;\\
  \ddot W_p(0)&=2\langle\dot c(0),\dot c(0)\rangle<0.
\end{align*}
因此存在 $\varepsilon>0$, 使对 $t\in(0,\varepsilon)$ 有 $W_p(t)<0$.

沿用与 Riemann 情形相同的思路 (参见第~\ref{chap:riemannian-manifolds}~章第~\ref{sec:minimizing-properties-of-geodesics}~节), 容易证明 $W_p$ 的水平曲面与过 $p$ 的测地线正交. 因而, 若 $c_v(t)=\exp_p(tv)$ 是初始条件为 $v\in T_pM$ 的测地线, 则有
\[
  (\operatorname{grad}W_p)_{c_v(1)}=a(v)\dot c_v(1),
\]
其中函数的梯度按 Riemann 情形中的方式定义 (不过要注意, 在 Lorentz 情形下, 若 $\operatorname{grad}f$ 为类时向量, 则光滑函数 $f$ 沿 $\operatorname{grad}f$ 的方向\textbf{减小}). 现在
\begin{align*}
  \left\langle(\operatorname{grad}W_p)_{c_v(t)},\dot c_v(t)\right\rangle
     &=\frac{\mathrm{d}}{\mathrm{d}t}W_p(c_v(t))
       =\frac{\mathrm{d}}{\mathrm{d}t}W_p(c_{tv}(1))\\
     &=\frac{\mathrm{d}}{\mathrm{d}t}\bigl(t^2W_p(c_v(1))\bigr)=2tW_p(c_v(1)),
\end{align*}
所以
\[
  \left\langle(\operatorname{grad}W_p)_{c_v(1)},\dot c_v(1)\right\rangle
     =2W_p(c_v(1)).
\]
另一方面,
\begin{align*}
  \left\langle(\operatorname{grad}W_p)_{c_v(1)},\dot c_v(1)\right\rangle
    &=\langle a(v)\dot c_v(1),\dot c_v(1)\rangle\\
    &=a(v)\langle v,v\rangle=a(v)W_p(c_v(1)).
\end{align*}
由此得出 $a(v)=2$, 因而
\[
  (\operatorname{grad}W_p)_{c_v(1)}=2\dot c_v(1).
\]
因此 $\operatorname{grad}W_p$ 与过 $p$ 的测地线相切, 且在朝向未来的测地线上指向未来.

假设 $W_p(t)<0$. 则 $(\operatorname{grad}W_p)_{c(t)}$ 为指向未来的类时向量, 所以
\[
  \dot W(t)=\left\langle(\operatorname{grad}W_p)_{c(t)},\dot c(t)\right\rangle<0,
\]
因为 $\dot c(t)$ 也是指向未来的类时向量 [参见 Exercise~\ref{ex:special-relativity}(\ref{ex:minkowski-orthogonality-and-time-orientation})]. 由此得出对所有 $t\in[0,1]$ 必有 $W_p(t)<0$. 特别地, $W_p(q)=W_p(1)<0$, 因此存在一条朝向未来的类时测地线连接 $p$ 至 $q$.

为证明论断 (\ref{cond:local-causal-future-is-chronological-closure}), 先证明 $\overline{I^+(p)}\subset J^+(p)$. 若 $q\in\overline{I^+(p)}$, 则 $q$ 是一列点 $q_n\in I^+(p)$ 的极限. 由 (\ref{cond:local-chronological-future-via-timelike-geodesics}), 有 $q_n=\exp_p(v_n)$, 其中 $v_n\in T_pM$ 是指向未来的类时向量. 由于 $\exp_p$ 是微分同胚, $v_n$ 收敛到一个指向未来的因果向量 $v\in T_pM$, 因而可由朝向未来的因果测地线从 $p$ 到达 $q=\exp_p(v)$. 反过来的包含关系 $J^+(p)\subset\overline{I^+(p)}$ 在一般情形下成立 (参见 Proposition~\ref{prop:global-chronological-and-causal-futures}).

最后, 由 (\ref{cond:local-causal-future-is-chronological-closure}) 和 $\exp_p$ 是到 $V$ 的微分同胚这一事实, (\ref{cond:local-boundary-via-null-geodesics}) 显然成立.
\end{RiemannProof}

利用这一局部性质, 可以证明如下整体结果.

\begin{RiemannProposition}\label{prop:global-chronological-and-causal-futures}
设 $(M,g)$ 为已作时间定向的时空, $p\in M$. 则:
\begin{enumerate}[label=(\arabic*),leftmargin=*]
\item $I^+(p)$ 是开集;
\item $J^+(p)\subset\overline{I^+(p)}$;
\item $I^+(p)=\operatorname{int}J^+(p)$;
\item 若 $r\in J^+(p)$ 且 $q\in I^+(r)$, 则 $q\in I^+(p)$;
\item 若 $r\in I^+(p)$ 且 $q\in J^+(r)$, 则 $q\in I^+(p)$.
\end{enumerate}
\end{RiemannProposition}

\begin{RiemannProof}
Exercise~\ref{ex:causality}(\ref{ex:properties-of-chronological-and-causal-futures}).
\end{RiemannProof}

广义双生子佯谬 [参见 Exercise~\ref{ex:special-relativity}(\ref{ex:generalized-twin-paradox-and-reverse-triangle})] 对一般时空也局部成立. 更确切地说, 有下列论断:

\begin{RiemannProposition}\label{prop:timelike-geodesic-maximizes-proper-time-locally}
设 $(M,g)$ 为已作时间定向的时空, $p_0\in M$, $V\subset M$ 为 $p_0$ 的测地凸开邻域. 将 $g$ 限制到 $V$ 所得的时空 $(V,g)$ 有以下性质: 若 $p,q\in V$ 且 $q\in I^+(p)$, $c$ 是连接 $p$ 至 $q$ 的类时测地线, $\gamma$ 是连接 $p$ 至 $q$ 的任意类时曲线, 则 $\tau(\gamma)\leq\tau(c)$, 当且仅当 $\gamma$ 是 $c$ 的重新参数化时取等号.
\end{RiemannProposition}

\begin{RiemannProof}
满足 $\gamma(0)=p$ 和 $\gamma(1)=q$ 的任意类时曲线 $\gamma:[0,1]\to V$ 都可对 $t\in[0,1]$ 写为
\[
  \gamma(t)=\exp_p(r(t)n(t)),
\]
其中 $r(t)\geq0$ 且 $\langle n(t),n(t)\rangle=-1$. 有
\[
  \dot\gamma(t)=(\exp_p)_*\bigl(\dot r(t)n(t)+r(t)\dot n(t)\bigr).
\]
由于 $\langle n(t),n(t)\rangle=-1$, 有 $\langle\dot n(t),n(t)\rangle=0$, 因而 $\dot n(t)$ 与函数 $v\mapsto\langle v,v\rangle$ 的水平曲面相切. 于是得到
\[
  \dot\gamma(t)=\dot r(t)X_{\gamma(t)}+Y(t),
\]
其中 $X$ 是过 $p$ 的类时测地线的单位切向量场, $Y(t)=r(t)(\exp_p)_*\dot n(t)$ 与 $W_p$ 的水平曲面相切 (因而与 $X_{\gamma(t)}$ 正交). 因此
\begin{align*}
  \tau(\gamma)
  &=\int_0^1\left|\left\langle
      \dot r(t)X_{\gamma(t)}+Y(t),\dot r(t)X_{\gamma(t)}+Y(t)
      \right\rangle\right|^{1/2}\,\mathrm{d}t\\
  &=\int_0^1\bigl(\dot r(t)^2-|Y(t)|^2\bigr)^{1/2}\,\mathrm{d}t\\
  &\leq\int_0^1\dot r(t)\,\mathrm{d}t=r(1)=\tau(c),
\end{align*}
这里用到了 $\gamma$ 为类时曲线, 对所有 $t\in[0,1]$ 有 $\dot r(t)>0$ (因为 $\dot\gamma$ 指向未来), 且 $\tau(c)=r(1)$ (因为 $q=\exp_p(r(1)n(1))$). 显然 $\tau(\gamma)=\tau(c)$ 当且仅当对所有 $t\in[0,1]$, $|Y(t)|\equiv0\iff Y(t)\equiv0$ ($Y(t)$ 是类空向量或零向量), 从而 $n$ 恒定. 在这种情形下, 除去重新参数化, $\gamma(t)=\exp_p(r(t)n)$ 就是过 $p$ 且初始条件为 $n\in T_pM$ 的测地线.
\end{RiemannProof}

零模测地线也有一条局部的刻画性质.

\begin{RiemannProposition}\label{prop:uniqueness-of-causal-curves-on-local-light-cone}
设 $(M,g)$ 为已作时间定向的时空, $p_0\in M$, $V\subset M$ 为 $p_0$ 的测地凸开邻域. 将 $g$ 限制到 $V$ 所得时空 $(V,g)$ 有以下性质: 若对 $p,q\in V$ 存在一条朝向未来的零模测地线 $c$ 连接 $p$ 至 $q$, 且 $\gamma$ 是连接 $p$ 至 $q$ 的因果曲线, 则 $\gamma$ 是 $c$ 的重新参数化.
\end{RiemannProposition}

\begin{RiemannProof}
由于 $p$ 与 $q$ 由零模测地线连接, 由 Proposition~\ref{prop:local-causal-structure-in-convex-normal-neighborhood} 得 $q\in J^+(p)\setminus I^+(p)$. 设 $\gamma:[0,1]\to V$ 是连接 $p$ 至 $q$ 的因果曲线. 则对所有 $t\in[0,1]$, 必有 $\gamma(t)\in J^+(p)\setminus I^+(p)$, 因为 $\gamma(t_0)\in I^+(p)$ 蕴含对所有 $t>t_0$ 有 $\gamma(t)\in I^+(p)$ (见 Proposition~\ref{prop:global-chronological-and-causal-futures}). 因而有
\[
  \left\langle(\operatorname{grad}W_p)_{\gamma(t)},\dot\gamma(t)\right\rangle=0,
\]
其中 $W_p$ 定义在 Proposition~\ref{prop:local-causal-structure-in-convex-normal-neighborhood} 的证明中. 对初始条件为 $v\in T_pM$ 的类时测地线 $c_v$ 已证明的公式
\[
  (\operatorname{grad}W_p)_{c_v(1)}=2\dot c_v(1)
\]
也必须对零模测地线成立 (由连续性). 因此, $\operatorname{grad}W_p$ 与织成 $J^+(p)\setminus I^+(p)$ 的零模测地线相切, 且指向未来. 由于 $\dot\gamma(t)$ 也指向未来, 得出 $\dot\gamma$ 与 $\operatorname{grad}W_p$ 成比例 [参见 Exercise~\ref{ex:special-relativity}(\ref{ex:generalized-twin-paradox-and-reverse-triangle})], 因而 $\gamma$ 必是一条零模测地线 (它必是 $c$) 的重新参数化.
\end{RiemannProof}

\begin{RiemannCorollary}\label{cor:global-boundary-causal-curves-are-null-geodesics}
设 $(M,g)$ 为已作时间定向的时空, 且 $p\in M$. 若 $q\in J^+(p)\setminus I^+(p)$, 则连接 $p$ 至 $q$ 的任意朝向未来的因果曲线都必为重新参数化的零模测地线.
\end{RiemannCorollary}

为作物理应用, 必须要求时空满足合理的因果性条件. 其中最简单的条件排除了时间旅行, 即排除了粒子返回自身过去历程中某个事件的可能性.

\begin{RiemannDefinition}\label{def:chronology-condition}
若时空 $(M,g)$ 不含闭类时曲线, 则称其满足\textbf{时序条件}.
\end{RiemannDefinition}

紧时空不满足这一条件:

\begin{RiemannProposition}\label{prop:compact-spacetime-has-closed-timelike-curves}
任意紧时空 $(M,g)$ 都含有闭类时曲线.
\end{RiemannProposition}

\begin{RiemannProof}
必要时取时间可定向的二重覆盖 [参见 Exercise~\ref{ex:causality}(\ref{ex:time-orientable-double-covering})], 可以假设 $(M,g)$ 已作时间定向. 由于对任意 $p\in M$, $I^+(p)$ 都是开集, 显然 $\{I^+(p)\}_{p\in M}$ 是 $M$ 的开覆盖. 若 $M$ 紧, 则可取有限子覆盖 $\{I^+(p_1),\ldots,I^+(p_N)\}$. 若对某个 $i\ne1$, 有 $p_1\in I^+(p_i)$, 则 $I^+(p_1)\subset I^+(p_i)$, 可以从子覆盖中删去 $I^+(p_1)$. 因此, 不失一般性可设 $p_1\in I^+(p_1)$, 从而存在以 $p_1$ 为起点和终点的闭类时曲线.
\end{RiemannProof}

对时空的因果行为可施加如下更强的限制:

\begin{RiemannDefinition}\label{def:stable-causality}
若时空 $(M,g)$ 存在\textbf{全局时间函数}, 即存在使 $\operatorname{grad}(t)$ 为类时向量场的光滑函数 $t:M\to\mathbb{R}$, 则称它是\textbf{稳定因果的}.
\end{RiemannDefinition}

特别地, 稳定因果时空是时间可定向的. 选取由 $-\operatorname{grad}(t)$ 定义的时间定向, 使 $t$ 沿朝向未来的类时曲线增大. 注意, 这意味着不可能存在闭类时曲线, 即任何稳定因果时空都满足时序条件. 事实上, 对稳定因果时空的任意小扰动仍满足时序条件 [参见 Exercise~\ref{ex:causality}(\ref{ex:small-perturbations-of-stably-causal-spacetime})].

设 $(M,g)$ 为已作时间定向的时空. 若 $\lim_{t\to b}c(t)$ 不存在, 则称光滑的朝向未来的因果曲线 $c:(a,b)\to M$ (可有 $a=-\infty$ 或 $b=+\infty$) \textbf{向未来不可延拓}. \textbf{向过去不可延拓}的因果曲线的定义类似. $S\subset M$ 的\textbf{未来依赖域} $D^+(S)$ 是满足如下条件的所有事件 $p\in M$ 所成的集合: 以 $p$ 为起点的任意向过去不可延拓的因果曲线都与 $S$ 相交. 因此, 对事件 $p\in D^+(S)$ 的任何因果影响一定曾在 $S$ 中某处留下记录, 可以预期由 $S$ 上的数据预测 $p$ 处所发生的事. 类似地, $S$ 的\textbf{过去依赖域} $D^-(S)$ 是满足如下条件的所有事件 $p\in M$ 所成的集合: 以 $p$ 为起点的任意向未来不可延拓的因果曲线都与 $S$ 相交. 因而事件 $p\in D^-(S)$ 的任何因果影响都会在 $S$ 中某处留下记录, 可以预期由 $S$ 上的数据反推 $p$ 处发生过的事. $S$ 的\textbf{依赖域}就是集合 $D(S)=D^+(S)\cup D^-(S)$.

设 $(M,g)$ 是具有时间函数 $t:M\to\mathbb{R}$ 的稳定因果时空. 若 $D(S_a)=M$, 则称水平集 $S_a=t^{-1}(a)$ 为 \textbf{Cauchy 超曲面}. 满足这一条件的时空具有特别好的因果性质.

\begin{RiemannDefinition}\label{def:globally-hyperbolic-spacetime}
若稳定因果时空具有一个时间函数, 其各个水平集都是 Cauchy 超曲面, 则称它是\textbf{整体双曲的}.
\end{RiemannDefinition}

注意, Cauchy 超曲面 $S_a$ 的未来和过去依赖域分别为 $D^+(S_a)=t^{-1}([a,+\infty))$ 和 $D^-(S_a)=t^{-1}((-\infty,a])$.

\begin{RiemannExercise}\label{ex:causality}
\begin{enumerate}[label=(\arabic*),ref=\arabic*,leftmargin=*]
\item\label{ex:time-orientable-double-covering}
(\emph{时间可定向的二重覆盖}) 使用与第~\ref{chap:differentiable-manifolds}~章的 Exercise~\ref{ex:orientability}(\ref{ex:orientable-double-covering}) 类似的思路, 证明若 $(M,g)$ 是时间不可定向的 Lorentz 流形, 则存在一个\textbf{时间可定向的二重覆盖}, 即时间可定向的 Lorentz 流形 $(\overline M,\overline g)$ 和局部等距映射 $\pi:\overline M\to M$, 使 $M$ 中每点在 $\pi$ 下有两个原像. 由此推出, 容许 Lorentz 度量的紧曲面只有环面 $T^2$ 和 Klein 瓶 $K^2$.
\item\label{ex:properties-of-chronological-and-causal-futures}
设 $(M,g)$ 为已作时间定向的时空, $p\in M$. 证明:
\begin{enumerate}[label=(\alph*),leftmargin=*]
\item $I^+(p)$ 是开集;
\item $J^+(p)$ 不一定是闭集;
\item $J^+(p)\subset\overline{I^+(p)}$;
\item $I^+(p)=\operatorname{int}J^+(p)$;
\item 若 $r\in J^+(p)$ 且 $q\in I^+(r)$, 则 $q\in I^+(p)$;
\item 若 $r\in I^+(p)$ 且 $q\in J^+(r)$, 则 $q\in I^+(p)$;
\item 可能有 $I^+(p)=M$.
\end{enumerate}
\item\label{ex:null-helical-curve-in-minkowski-spacetime}
考虑三维 Minkowski 时空 $(\mathbb{R}^3,g)$, 其中
\[
  g=-\mathrm{d}t\otimes\mathrm{d}t+\mathrm{d}x\otimes\mathrm{d}x
    +\mathrm{d}y\otimes\mathrm{d}y.
\]
设 $c:\mathbb{R}\to\mathbb{R}^3$ 为曲线 $c(t)=(t,\cos t,\sin t)$. 证明尽管对所有 $t\in\mathbb{R}$, $\dot c(t)$ 都是零模向量, 但对所有 $t>0$ 都有 $c(t)\in I^+(c(0))$. 此曲线表示什么类型的运动?
\item\label{ex:small-perturbations-of-stably-causal-spacetime}
设 $(M,g)$ 为稳定因果时空, $h$ 为任意具有紧支集的对称 $(2,0)$-张量场. 证明对于足够小的 $|\varepsilon|$, 张量场 $g_\varepsilon:=g+\varepsilon h$ 仍是 $M$ 上的 Lorentz 度量, 且 $(M,g_\varepsilon)$ 满足时序条件.
\item\label{ex:chronological-but-not-stably-causal-quotient}
设 $(M,g)$ 是二维 Minkowski 时空对由映射 $f(t,x)=(t+1,x+1)$ 生成的离散等距映射群取商所得. 证明 $(M,g)$ 满足时序条件, 但存在 $(M,g)$ 的任意小扰动 [按 Exercise~\ref{ex:causality}(\ref{ex:small-perturbations-of-stably-causal-spacetime}) 的意义], 使得时序条件不再满足.
\item\label{ex:future-domain-of-dependence-examples}
设 $(M,g)$ 为已作时间定向的时空, $S\subset M$. 证明:
\begin{enumerate}[label=(\alph*),leftmargin=*]
\item $S\subset D^+(S)$;
\item $D^+(S)$ 不一定是开集;
\item $D^+(S)$ 不一定是闭集.
\end{enumerate}
\item\label{ex:slit-minkowski-spacetime}
设 $(M,g)$ 是从二维 Minkowski 时空中去掉 $x$ 轴正半轴所得的二维时空 (参见 Fig.~\ref{fig:stably-causal-not-globally-hyperbolic-spacetime}). 证明:
\begin{enumerate}[label=(\alph*),leftmargin=*]
\item $(M,g)$ 稳定因果, 但不整体双曲;
\item 存在 $p,q\in M$, 使 $J^+(p)\cap J^-(q)$ 不紧;
\item 存在 $p,q\in M$ 且 $q\in I^+(p)$, 使连接 $p$ 至 $q$ 的类时曲线长度的上确界, 不被任何类时曲线取到.
\end{enumerate}
\item\label{ex:global-hyperbolicity-of-product-spacetime}
设 $(\Sigma,h)$ 为三维 Riemann 流形. 证明时空 $(M,g)=(\mathbb{R}\times\Sigma,-\mathrm{d}t\otimes\mathrm{d}t+h)$ 整体双曲当且仅当 $(\Sigma,h)$ 完备.
\item\label{ex:globally-hyperbolic-spacetime-examples}
证明下列时空都是整体双曲的:
\begin{enumerate}[label=(\alph*),leftmargin=*]
\item Minkowski 时空;
\item FLRW 时空;
\item Schwarzschild 时空的区域 $\{r>2m\}$;
\item Schwarzschild 时空的区域 $\{r<2m\}$.
\end{enumerate}
\item\label{ex:globally-hyperbolic-spacetime-splitting}
设 $(M,g)$ 为具有 Cauchy 超曲面 $S$ 的整体双曲时空. 证明 $M$ 与 $\mathbb{R}\times S$ 微分同胚.
\end{enumerate}
\end{RiemannExercise}
\RGfigure[.66\textwidth]{stably-causal-not-globally-hyperbolic-spacetime}{稳定因果但不整体双曲的时空}
